% SOURCE_AUTHORITY: sga5_fr_workpass.tex, lines 3526--3618 (LF-normalized unit).
% SOURCE_SHA256: 791F4EFFC5E02832D5D77ED03518C8156D6F07E4C8238B03545DB93D883FBB28
\subsection*{5.2. Composição das correspondências}

Sejam, para $i=1,2,3$, $X_i$ um $S$-esquema, $X_{ij}=X_i\times_SX_j$, $X=X_1\times_SX_2\times_SX_3$, $p_i:X\to X_i$, $p_{ij}:X\to X_{ij}$ as projeções e $L_i\in\operatorname{ob}D_{\mathrm{ctf}}(X_i)$. Tem-se $X=X_{12}\times_{X_2}X_{23}$, e
\[
(DL_1\otimes_S^{\mathbf L}L_2)\otimes_{X_2}^{\mathbf L}(DL_2\otimes_S^{\mathbf L}L_3)
=p_1^*DL_1\otimes^{\mathbf L}p_2^*L_2\otimes^{\mathbf L}p_2^*DL_2\otimes^{\mathbf L}p_3^*L_3.
\]
A flecha de avaliação $L_2\otimes^{\mathbf L}DL_2\to K_{X_2}$ define, portanto, um emparelhamento
\begin{equation}
\tag{5.2.1}
(DL_1\otimes_S^{\mathbf L}L_2)\otimes_{X_2}^{\mathbf L}(DL_2\otimes_S^{\mathbf L}L_3)\longrightarrow p_{13}^!(DL_1\otimes_S^{\mathbf L}L_3),
\end{equation}
identificando-se o segundo membro com $(DL_1\otimes_S^{\mathbf L}L_3)\otimes_S^{\mathbf L}K_{X_2}$ mediante Künneth (1.7.3). Levando em conta 3.1.1 e 3.1.2, 5.2.1 pode ser reescrita na forma
\begin{equation}
\tag{5.2.2}
\uRHom_S(L_1,L_2)\otimes_{X_2}^{\mathbf L}\uRHom_S(L_2,L_3)\longrightarrow p_{13}^!\uRHom_S(L_1,L_3)=\uRHom(p_1^*L_1,p_3^!L_3),
\end{equation}
onde a identificação no segundo membro é dada pelo isomorfismo de indução (SGA~4~XVIII 3.1.12.2), daí se obtendo um emparelhamento, denominado \emph{composição de correspondências},
\begin{equation}
\tag{5.2.3}
\Hom_S(L_1,L_2)\otimes\Hom_S(L_2,L_3)\longrightarrow\Hom(p_1^*L_1,p_3^!L_3),
\end{equation}
que denotaremos por $(u,v)\mapsto vu$. Verifica-se que, se $p_{ij,i}:X_{ij}\to X_i$ e $p_{ij,j}:X_{ij}\to X_j$ denotam as projeções, 5.2.2 também pode ser obtida pela composição da flecha
\[
\uRHom_S(L_1,L_2)\otimes_{X_2}^{\mathbf L}\uRHom_S(L_2,L_3)
\longrightarrow
\uRHom(p_1^*L_1,p_{12}^*p_{12,2}^!L_2)
\otimes
\uRHom(p_{23}^!p_{23,2}^*L_2,p_3^!L_3)
\]
(produto tensorial das flechas canônicas $p_{12}^*\uRHom_S(L_1,L_2)\to\uRHom(p_1^*L_1,p_{12}^*p_{12,2}^!L_2)$ e $p_{23}^*\uRHom_S(L_2,L_3)\to\uRHom(p_{23}^!p_{23,2}^*L_2,p_3^!L_3)$), da flecha definida pela flecha de mudança de base $p_{12}^!p_{12,2}^*L_2\to p_{23}^!p_{23,2}^*L_2$, e de uma flecha de composição habitual nas $\uRHom$.

Sejam $c':C'\to X_{12}$ e $c'':C''\to X_{23}$ correspondências, e seja $c:C=C'\times_{X_2}C''\to X$ a correspondência ``composta''. De 5.2.2 deduz-se (ou define-se diretamente, mediante uma composição análoga à que acabamos de descrever) um emparelhamento
\begin{equation}
\tag{5.2.4}
\uRHom(c_1'{}^*L_1,c_2'{}^!L_2)\otimes_{X_2}^{\mathbf L}\uRHom(c_2''{}^*L_2,c_3''{}^!L_3)
\longrightarrow\uRHom(c_1^*L_1,c_3^!L_3),
\end{equation}
de onde se obtém um emparelhamento ``composição'' $(u,v)\mapsto vu$:
\begin{equation}
\tag{5.2.5}
\Hom(c_1'{}^*L_1,c_2'{}^!L_2)\otimes\Hom(c_2''{}^*L_2,c_3''{}^!L_3)
\longrightarrow\Hom(c_1^*L_1,c_3^!L_3).
\end{equation}

Verifica-se que a composição das correspondências é compatível com as imagens diretas, no sentido seguinte. Sejam, para $1\leq i\leq 3$, $f_i:X_i\to Y_i$ um $S$-morfismo, sendo $f_2$ próprio; sejam $c'$, $c''$, $c$ morfismos como antes, com $c'$ e $c''$ próprios, $d':D'\to Y_{12}$, $d'':D''\to Y_{23}$ morfismos próprios e $d:D=D'\times_{Y_2}D''\to Y$, e suponhamos dados quadrados comutativos
\[
\begin{tikzcd}[column sep=large,row sep=large]
X_{12} \arrow[d,"f_{12}"'] & C' \arrow[l] \arrow[d] \\
Y_{12} & D' \arrow[l]
\end{tikzcd}
\qquad
\begin{tikzcd}[column sep=large,row sep=large]
X_{23} \arrow[d,"f_{23}"'] & C'' \arrow[l] \arrow[d] \\
Y_{23} & D'' \arrow[l]
\end{tikzcd}
\]
Tem-se então um quadrado comutativo
\begin{equation}
\tag{5.2.6}
\begin{tikzcd}[column sep=large,row sep=large]
\Hom(c_1'{}^*L_1,c_2'{}^!L_2)\otimes\Hom(c_2''{}^*L_2,c_3''{}^!L_3) \arrow[r,"5.2.5"] \arrow[d,"f_{12*}\otimes f_{23*}"']
& \Hom(c_1^*L_1,c_3^!L_3) \arrow[d,"f_{13*}"] \\
\Hom(d_1'{}^*M_1,d_2'{}^!M_2)\otimes\Hom(d_2''{}^*M_2,d_3''{}^!M_3) \arrow[r,"5.2.5"]
& \Hom(d_1^*M_1,d_3^!M_3),
\end{tikzcd}
\end{equation}
onde $M_1=f_{1!}L_1$, $M_2=f_{2*}L_2$, $M_3=f_{3*}L_3$, e tem-se um quadrado comutativo análogo com $\uRHom$, cuja escrita deixamos ao leitor.

Por fim, suponhamos que $X_1=X_3$, $L_1=L_3$. O emparelhamento fundamental 4.1.2 pode ser interpretado como a composta do emparelhamento
\begin{equation}
\tag{5.2.7}
q_1^*DL_1\otimes^{\mathbf L}q_2^*L_2\otimes^{\mathbf L}q_2^*DL_2\otimes^{\mathbf L}q_1^*L_1
\longrightarrow
q_1^*DL_1\otimes^{\mathbf L}q_2^*K_{X_2}\otimes^{\mathbf L}q_1^*L_1
\end{equation}
definido por $L_2\otimes^{\mathbf L}DL_2\to K_{X_2}$, ou, o que é o mesmo, induzido por 5.2.1 sobre $X_{12}$ por meio da diagonal $X_{12}\to X_1\times_SX_2\times_SX_1$ (onde $q_i:X_{12}\to X_i$ denota a projeção), e do emparelhamento
\begin{equation}
\tag{5.2.8}
q_1^*DL_1\otimes^{\mathbf L}q_2^*K_{X_2}\otimes^{\mathbf L}q_1^*L_1\longrightarrow K_{X_{12}}
\end{equation}
definido por $DL_1\otimes^{\mathbf L}L_1\to K_{X_1}$. Daí se deduz que, para $u\in\Hom_S(L_1,L_2)$, $v\in\Hom_S(L_2,L_1)$, tem-se, em $H^0(X_1\times_SX_2,K_{X_1\times_SX_2})$,
\begin{equation}
\tag{5.2.9}
\langle u,v\rangle=\langle vu,1\rangle,
\end{equation}
onde $1$ designa a correspondência de identidade de $L_1$ a $L_1$ sobre $X_1\times_SX_1$, com suporte na diagonal (3.2 b)), e o produto cup do segundo membro é tomado no sentido de 4.2.5, tendo em conta que o produto fibrado $X\times_{(X_1\times_SX_1)}X_1$ (por meio de $p_{13}:X\to X_{13}=X_1\times_SX_1$ e da diagonal $X_1\to X_1\times_SX_1$) identifica-se com $X_1\times_SX_2$. Da mesma forma, se vê, mais geralmente, que, se $c':C'\to X_{12}$ e $c'':C''\to X_{21}$ são correspondências, $c:C=C'\times_{X_2}C''\to X_1\times_SX_2\times_SX_1$ é a correspondência composta e $d:C'\times_{X_{12}}C''=D\to X_{12}$ a correspondência de intersecção, então, para $u\in\Hom(c_1'{}^*L_1,c_2'{}^!L_2)$, $v\in\Hom(c_2''{}^*L_2,c_1''{}^!L_1)$, temos a seguinte igualdade em $H^0(D,K_D)$:
\begin{equation}
\tag{5.2.10}
\langle u,v\rangle=\langle vu,1\rangle,
\end{equation}
onde $1$ designa, como antes, a correspondência diagonal, e $D$ identifica-se com o produto fibrado de $C\to X_1\times_SX_1$ e a diagonal $X_1\to X_1\times_SX_1$. As fórmulas 5.2.9 e 5.2.10 generalizam uma fórmula bem conhecida no caso dos coeficientes constantes ([12] 1.3.6 (ii) c)).
